\section{模型的建立与数值方法}

\subsection{建模准备与总体框架}

四个问题研究的是同一个对象，即圆柱形药材在热风中的升温与失水，控制方程、边界条件与初值完全一致。
因此我们不逐问重复建模，而是先建立一套统一的径向轴对称模型，四个问题都在它上面求解，
只在物性经验式、环境时程与计算区域三处分支：问题一取问题给定的常物性，
问题二与问题三取附录 3 的温度—浓度双向耦合经验式，问题四再叠加由附件 2 决定的尺寸收缩。
\\这样处理既免去重复推导，也使四问的结果具有直接可比性。本节先给出控制方程与定解条件，
再给出统一的数值方法，各问的具体参数与求解结果分别见第六章至第九章。

\subsection{控制方程的建立}

\subsubsection{由守恒律导出控制方程}

药材内部既无内热源也无化学反应，故热量只靠导热传递，水分只由浓度梯度驱动。
由微元法，在药材内取半径由 $r$ 到 $r+\Delta r$ 的一层薄环，并沿轴向取单位长度，由\textbf{能量守恒定律}环内热量的变化率应等于
两侧界面净导入的导热热流，即

\begin{equation}
\rho c_p\frac{\partial T}{\partial t}\cdot 2\pi r\,\Delta r
=\left.\left(k\cdot 2\pi r\frac{\partial T}{\partial r}\right)\right|_{r+\Delta r}
-\left.\left(k\cdot 2\pi r\frac{\partial T}{\partial r}\right)\right|_{r}.
\label{eq:energybalance}
\end{equation}

两端同除以 $2\pi r\,\Delta r$，再令 $\Delta r\to 0$，右端即成为散度形式，于是得到能量方程

\begin{equation}
\rho c_p\frac{\partial T}{\partial t}
=\frac{1}{r}\frac{\partial}{\partial r}\left(k\,r\,\frac{\partial T}{\partial r}\right).
\label{eq:energy}
\end{equation}

对水分做同样的守恒分析：环内水量的变化率等于两侧界面净导入的水分质量流，
而质量流的本构关系即\textbf{ Fick 定律}，即通量等于扩散系数乘以浓度梯度的负值，于是

\begin{equation}
\frac{\partial C}{\partial t}
=\frac{1}{r}\frac{\partial}{\partial r}\left(D\,r\,\frac{\partial C}{\partial r}\right).
\label{eq:mass}
\end{equation}

观察到，问题一中物性 $k$、$D$ 与状态无关，式~\eqref{eq:energy} 与式~\eqref{eq:mass} 彼此独立，
温度场与水分场可以分别求解。

需要说明的是，式~\eqref{eq:mass} 中不含密度：本文按题目明文的 Fick 标准形式建立质量方程，
密度、比热容与导热系数只出现在能量方程中。这一读法决定了后文全部时长结论的适用范围，
其影响量级在灵敏度分析中定量给出。

\subsubsection{边界条件与初值}

\begin{tabular}{@{}p{1.4em}p{0.85\textwidth}@{}}
① & 在中心 \(r=0\) 处，由轴对称性给出零梯度条件；\\[2pt]
② & 在表面 \(r=R\) 处，热流与质量流均由对流控制。\\
\end{tabular}

把外法向通量写开，即得第三类（Robin）边界条件
\begin{equation}
\left.-k\frac{\partial T}{\partial r}\right|_{r=R}
=h\left(T(R,t)-T_\infty(t)\right),
\label{eq:bc}
\end{equation}

\begin{equation}
\left.-D\frac{\partial C}{\partial r}\right|_{r=R}
=h_m\left(C(R,t)-C_\infty(t)\right).
\label{eq:bcm}
\end{equation}

其中 $h$ 与 $h_m$ 分别为对流传热系数与对流传质系数，$T_\infty$ 与 $C_\infty$ 为烘房环境的
温度与水分浓度。初始时刻温度与含水率在全域均匀，即 $T(r,0)=T_0$、$C(r,0)=C_0$。\\
[2pt]环境序列由附件 1 的实测数据给出（$0$--$14400$~s，步长 $60$~s），采样点之间的时刻按相邻两点的分段线性插值取值。
$14400$~s 之后记录终止，本文取末段（约 $t\ge9000$~s）稳定平台的实测均值并保持不变，
即 $T_\infty=50$~$^\circ$C、$C_\infty=0.05$~kg/kg，以支撑问题二至问题四 的计算。

\subsubsection{特征时间与量级分析}

数值求解之前，先由控制方程判断两个阶段的主控机制，为时间步长与验证设计定标。
对式~\eqref{eq:mass} 与式~\eqref{eq:bcm} 作无量纲化，取 $\xi=r/R$、$u=(C-C_\infty)/(C_0-C_\infty)$，得

\begin{equation}
\frac{\partial u}{\partial\tau}
=\frac{1}{\xi}\frac{\partial}{\partial\xi}\left(\xi\frac{\partial u}{\partial\xi}\right),\qquad
\tau=\frac{Dt}{R^2},
\label{eq:massnd}
\end{equation}

\begin{equation}
\left.-\frac{\partial u}{\partial\xi}\right|_{\xi=1}
=\text{Bi}_m\,u(1,\tau),\qquad
\text{Bi}_m=\frac{h_mR}{D}=\frac{R/D}{1/h_m}.
\label{eq:bcnd}
\end{equation}

式~\eqref{eq:massnd} 不含任何参数，物性只由式~\eqref{eq:bcnd} 进入，
可见水分过程仅由 Fourier 数 $\tau$ 与传质 Biot 数 $\text{Bi}_m$ 支配。

$\tau$ 表示已过去的时间相当于多少个特征扩散时间 $R^2/D$。预热阶段处于湿端
（$C\approx1.5$--$2.5$~kg/kg、$T\approx28$~$^\circ$C），$R^2/D\approx20$--$22$~h，
而该阶段只有 $30$~min，$t=1800$~s 时 $\tau\approx0.025$、穿透深度 $\sqrt{Dt}\approx0.3$~cm，
不足半径的两成，故\textbf{失水只能发生在表层}。
\\$\text{Bi}_m$ 是内部扩散阻力与外部对流传质阻力之比，本题湿端约 $3$、干端约 $20$，全程大于 $1$，
说明\textbf{表面传质阻力小于内部扩散阻力，干燥由内部扩散控制}。

这两个量级判断将在第六、七章的灵敏度分析中被定量证实。

\subsection{数学求解过程}

四个问题共用同一条求解链条：把空间导数离散成守恒通量形式，把时间导数按 $\theta$ 加权格式推进，
问题四再叠加移动边界的表观对流项，每个时间步归结为一个三对角方程组并用追赶法求解，
最后在题目要求的报告时刻与位置取值。四问的差别只在代入哪套物性经验式、
是否启用坐标变换以及计算区域如何取，链条本身不变。完整流程见图~\ref{fig:flow}。

\begin{figure}[H]
  \centering
  \includegraphics[width=0.94\textwidth]{fig07_solution_flow.pdf}
  \caption{统一求解流程。由环境序列与几何数据出发，经控制体权重与界面系数构造，
  按问题选取物性经验式与时间格式，最后由追赶法逐时步推进并输出报告值。}
  \label{fig:flow}
\end{figure}

下面按空间离散、时间离散、移动边界处理与验证设计的顺序给出各环节。

\subsubsection{空间离散：对偶控制体的顶点有限体积格式}

空间离散采用有限体积法：把守恒律写在每个小控制体上，界面通量一进一出相互抵消，
因此水分不会凭空产生或消失。下面按“取控制体—写通量”两步构造格式。

把区间 $[0,R]$ 等分为 $N$ 段，节点为 $r_j=j\Delta r$。以每个节点为中心取对偶控制体：
内部节点的控制体为 $[r_j-\Delta r/2,\;r_j+\Delta r/2]$，首尾节点各取半格。
把式~\eqref{eq:mass} 在每个控制体上作径向加权积分（权为 $r$，即环面积因子），得

\begin{equation}
\int_{V_j} r\frac{\partial C}{\partial t}\,dr
=\left[\,r D\frac{\partial C}{\partial r}\right]_{\text{两侧界面}},
\label{eq:fv}
\end{equation}

右端只出现界面上的通量，这正是有限体积法的出发点。把左端的积分近似为
$w_j\,\mathrm{d}C_j/\mathrm{d}t$，其中权重 $w_j=\int_{V_j}r\,\mathrm{d}r$ 为对偶控制体上径向坐标的积分，
把右端的界面通量用相邻节点的差商近似，即得半离散方程

\begin{equation}
w_j\frac{\mathrm{d}C_j}{\mathrm{d}t}
=r_{j+\frac12}D_{j+\frac12}\frac{C_{j+1}-C_j}{\Delta r}
-r_{j-\frac12}D_{j-\frac12}\frac{C_j-C_{j-1}}{\Delta r}.
\label{eq:semi}
\end{equation}

图~\ref{fig:mesh} 给出这一离散的几何含义：左侧在圆柱截面上取一行径向节点，
右侧放大后可以看到节点、对偶控制体与界面通量的相对位置。
本文取 $R=2.0$~cm、$N=800$ 个节点，故 $\Delta r=R/N=25$~$\mu$m。

\begin{figure}[H]
  \centering
  \includegraphics[width=0.98\textwidth]{fig00_mesh_schematic.pdf}
  \caption{顶点型有限体积离散示意}
  \label{fig:mesh}
\end{figure}

\subsubsection{两个端点的处理}

在中心 $r=0$ 处，对偶控制体的内侧界面半径为 $0$，式~\eqref{eq:semi} 中对应的通量项自动为零，
轴对称条件\textbf{由格式本身满足}，无需对 $1/r$ 的奇点作任何特殊处理。
在表面 $r=R$ 处，节点本身就是边界点，把式~\eqref{eq:bcm} 的通量代入式~\eqref{eq:fv} 的界面项，得

\begin{equation}
w_N\frac{\mathrm{d}C_N}{\mathrm{d}t}
=r_{N-\frac12}D_{N-\frac12}\frac{C_{N-1}-C_N}{\Delta r}
-R\,h_m\left(C_N-C_\infty\right).
\label{eq:surf}
\end{equation}

式~\eqref{eq:surf} 说明：报告的表面值就是解变量 $C_N$ 本身，不需要由最外层单元中心回算。
这避免了半单元重构引入的一阶误差，也正是我们选择顶点格式作为主要求解格式的原因。

\subsubsection{时间离散：$\theta$ 加权与启动步}

对半离散系统 $\mathbf{w}\,\mathrm{d}\mathbf{y}/\mathrm{d}t=\mathbf{Q}(\mathbf{y})$ 取 $\theta$ 加权，得

\begin{equation}
\mathbf{y}^{n+1}-\mathbf{y}^{n}
=\Delta t\left[\theta\,\mathbf{Q}^{n+1}+(1-\theta)\mathbf{Q}^{n}\right].
\label{eq:theta}
\end{equation}

当 $\theta=1$ 时为后向 Euler，精度一阶但耗散强、无条件稳定；当 $\theta=1/2$ 时为
Crank--Nicolson，精度二阶且 A 稳定。本文取 $\theta=1/2$，但\textbf{前 $3$ 步改用 $\theta=1$}：
初值处的温湿场与边界驱动之间存在突变，纯 Crank--Nicolson 会把这一突变转化为随时间衰减的振荡，
以后向 Euler 启动即可将其抹平；三步之后解已充分光滑，再切回二阶格式不影响整体精度。

由于物性依赖温度与含水率，每一步的系数取上一时间层的值，使每个时间步仍是一个形如

\begin{equation}
a_jy_{j-1}+b_jy_j+c_jy_{j+1}=d_j
\label{eq:tridiag}
\end{equation}

的\textbf{三对角方程组}，可用追赶法在 $O(N)$ 时间内求解。系数滞后带来的时间一阶误差
由时间加密实验核验（见第七、八章）。

\subsubsection{移动边界的坐标变换}

问题四中半径随时间收缩为 $R(t)$，固定网格无法贴合移动的边界。引入归一化坐标 $\xi=r/R(t)$，
把未知的移动域 $[0,R(t)]$ 映射为固定的 $[0,1]$，这一步把“网格跟着边界走”转化为“边界不动、
方程多一项”。设 $C(r,t)=\tilde C(\xi,t)$，由链式法则

\begin{equation}
\left.\frac{\partial C}{\partial t}\right|_r
=\left.\frac{\partial \tilde C}{\partial t}\right|_\xi
+\frac{\partial \tilde C}{\partial \xi}\cdot\left.\frac{\partial \xi}{\partial t}\right|_r,
\qquad \left.\frac{\partial \xi}{\partial t}\right|_r=-\frac{r\dot R}{R^2}=-\frac{\xi\dot R}{R},
\label{eq:chain}
\end{equation}

同时 $\partial/\partial r=R^{-1}\,\partial/\partial\xi$，代入式~\eqref{eq:mass} 即得

\begin{equation}
\left.\frac{\partial \tilde C}{\partial t}\right|_\xi
=\frac{1}{\xi R^2}\frac{\partial}{\partial \xi}\left(D\,\xi\,\frac{\partial \tilde C}{\partial \xi}\right)
+\underbrace{\frac{\xi\dot R}{R}\frac{\partial \tilde C}{\partial \xi}}_{\text{表观对流项}} .
\label{eq:moving}
\end{equation}

这一项的物理含义是清楚的：干基含水率按定义不随几何收缩而改变，因此在固定于材料的归一化坐标下，
收缩表现为一个把浓度剖面按 $\xi\dot R/R$ 平流的对流项，而不是一个凭空产生或消灭水分的源汇。
$\dot R$ 由附件 2 的单调插值解析求导得到，是已知函数，因此式~\eqref{eq:moving} 仍是线性方程组，
无需重构网格即可求解。

\subsubsection{移动边界的守恒性核验与不变量检验}

移动边界下的守恒性不易目测，本文把它写成一个可以直接核验的离散恒等式。
把式~\eqref{eq:moving} 两端同乘 $\xi$ 并在 $[0,1]$ 上积分，左端为
$\frac{\mathrm{d}}{\mathrm{d}t}\int_0^1\xi\tilde C\,\mathrm{d}\xi$（$\xi$ 域固定，可移入导数）。
右端第一项按分部积分给出界面项 $\left[\xi D\,\partial_\xi\tilde C/R^2\right]_0^1
=-\frac{h_m}{R}(C_N-C_\infty)$，第二项给出
$\frac{\dot R}{R}\int_0^1\xi^2\partial_\xi\tilde C\,\mathrm{d}\xi
=\frac{\dot R}{R}\left(C_N-2\int_0^1\xi\tilde C\,\mathrm{d}\xi\right)$。
把两个积分用归一化坐标下对偶控制体的权重 $\omega_j=\int_{\text{cell}}\xi\,\mathrm{d}\xi$ 离散
（它与式~\eqref{eq:fv} 中的 $w_j$ 相差常数因子 $R^2$），整理得

\begin{equation}
\frac{\mathrm{d}}{\mathrm{d}t}\sum_j \omega_jC_j
=-2\frac{\dot R}{R}\sum_j \omega_jC_j+\frac{\dot R}{R}C_N-\frac{h_m}{R}\left(C_N-C_\infty\right).
\label{eq:conservation}
\end{equation}

式~\eqref{eq:conservation} 是一个\textbf{可以直接逐时步核验的离散恒等式}：
左端由数值解给出，右端由当步的系数给出，二者之差即格式的守恒性残差。

再设一条更锋利的判据：令 $h_m=0$ 但保持收缩，则式~\eqref{eq:moving} 退化为纯平流方程，
其解析解要求 $\tilde C(\xi,t)\equiv C_0$ 恒成立——因为干基含水率不随几何缩小而改变。
若表观对流项的符号或系数写错，收缩就会凭空产生或消灭水分，该检验必然立刻暴露。
这里我们把这条判据称为\textbf{解析不变量检验}，它是针对移动边界项最有针对性的一条验证设计，
其结果见第九章。

\subsubsection{求解流程与统一验证设计}

综合以上推导，四个问题共用同一套求解流程：读取环境序列与几何数据；构造对偶控制体权重与界面系数；
按问题选择物性经验式与时间格式；组装三对角方程组并用追赶法逐时步推进；在报告时刻按题目要求的
位置取值输出。其中问题四额外按 $\dot R$ 叠加表观对流项。完整流程见图~\ref{fig:flow}。

每一问都配齐三条相互独立的实现路径：本文方法（顶点有限体积）、对照方法（格心有限体积，
其表面值由式~\eqref{eq:bcm} 回算）与独立复算（常物性下的级数解析解，以及式~\eqref{eq:conservation}
给出的守恒恒等式与本节的不变量检验）。三条路径不共享数值输入，只在其结果上相互比对，
以保证“算对了”这一判断不是自证。各问的比对结果分别见第六至第九章。
